\documentclass[12pt,a4paper]{report}

% =========================
% PACKAGE
% =========================
\usepackage[utf8]{inputenc}
\usepackage[vietnamese]{babel}
\usepackage{graphicx}
\usepackage{amsmath}
\usepackage{geometry}
\usepackage{setspace}
\usepackage{hyperref}

% =========================
% PAGE SETUP
% =========================
\geometry{
	left=3cm,
	right=2.5cm,
	top=1cm,
	bottom=2.5cm
}

\onehalfspacing

% =========================
% BEGIN DOCUMENT
% =========================
\begin{document}
	
	% =========================
	% COVER PAGE
	% =========================
	\begin{titlepage}
		\centering
		
		\vspace*{1cm}
		
		{\Large \textbf{TRƯỜNG ĐẠI HỌC GIAO THÔNG VẬN TẢI THÀNH PHỐ HỒ CHÍ MINH}}\\[0.5cm]
		
		{\large KHOA CÔNG NGHỆ THÔNG TIN}\\[1cm]
		
		{\Huge \textbf{BÁO CÁO ĐỒ ÁN}}\\[1cm]
		
		{\LARGE
			\textbf{
				XÂY DỰNG CHATBOT HỎI ĐÁP\\
				DỰA TRÊN TÀI LIỆU MÔN HỌC\\
				BẰNG TIẾNG VIỆT
		}}\\[1cm]
		
		{\Large
			So sánh hiệu quả giữa RAG và Fine-tuning
		}
		
		\vfill
		
		\begin{flushleft}
			\textbf{Nhóm thực hiện:} Nhóm 9\\
			\textbf{Giảng viên hướng dẫn:} Ths:Nguyễn Văn Chiến\\
			\textbf{Môn học:} Công nghệ phần mềm\\
			\textbf{Thành viên nhóm:} \\
			                 Đào Văn Khánh\\
		                    Nguyễn Ninh Công Danh\\
		                  	Phan Liên Hưng\\
		                  	Trần Thị Kiều Trinh\\
		                  	Nguyễn Hiển Đạt\\
		                  	Lương Thành Đạt\\
		                  	Phạm Hoàng Dũng\\
		                 
		                  
		\end{flushleft}
		
		\vspace{2cm}
		
		{\large \today}
		
	\end{titlepage}
	
	% =========================
	% TABLE OF CONTENTS
	% =========================
	\tableofcontents
	\newpage
	
	% =========================
	% CHAPTER 1
	% =========================
	\chapter{Giới thiệu đề tài}
	
	\section{Lý do chọn đề tài}
	
	Sự phát triển của trí tuệ nhân tạo và các mô hình ngôn ngữ lớn đã mở ra nhiều ứng dụng trong giáo dục và hỗ trợ học tập. Trong đó, chatbot hỏi đáp tài liệu học tập là một hướng nghiên cứu được quan tâm nhằm hỗ trợ sinh viên tiếp cận kiến thức hiệu quả hơn.
	
	\section{Mục tiêu đề tài}
	
	Đề tài hướng đến xây dựng chatbot hỗ trợ sinh viên hỏi đáp dựa trên tài liệu môn học bằng tiếng Việt, đồng thời nghiên cứu và so sánh hiệu quả giữa hai phương pháp RAG và Fine-tuning.
	
	\section{Phạm vi đề tài}
	
	Hệ thống hỗ trợ:
	\begin{itemize}
		\item Upload tài liệu PDF/DOCX/PPT
		\item Chatbot hỏi đáp theo tài liệu
		\item Trích dẫn nguồn tài liệu
		\item Benchmark RAG và Fine-tuning
	\end{itemize}
	
	% =========================
	% CHAPTER 2
	% =========================
	\chapter{Cơ sở lý thuyết}
	
	\section{Retrieval-Augmented Generation}
	
	RAG là phương pháp kết hợp truy xuất dữ liệu và mô hình sinh ngôn ngữ để cải thiện chất lượng trả lời của chatbot.
	
	\section{Fine-tuning}
	
	Fine-tuning là quá trình huấn luyện lại mô hình AI trên dữ liệu chuyên biệt để tăng độ chính xác cho bài toán cụ thể.
	
	\section{Embedding Model}
	
	Các embedding model được sử dụng:
	\begin{itemize}
		\item multilingual-e5-base
		\item bge-m3
		\item PhoBERT-base
		\item text-embedding-3-small
	\end{itemize}
	
	% =========================
	% CHAPTER 3
	% =========================
	\chapter{Phân tích và thiết kế hệ thống}
	
	\section{Use Case Diagram}
	
	Mô tả các chức năng chính của hệ thống.
	
	\section{ERD}
	
	Thiết kế cơ sở dữ liệu và các mối quan hệ giữa các thực thể.
	
	\section{Activity Diagram}
	
	Mô tả luồng hoạt động của hệ thống chatbot.
	
	% =========================
	% CHAPTER 4
	% =========================
	\chapter{Kiến trúc hệ thống}
	
	\section{Pipeline hệ thống}
	
	Pipeline của hệ thống bao gồm:
	\begin{enumerate}
		\item Upload tài liệu
		\item Chunking
		\item Embedding
		\item Vector Database
		\item Retrieval
		\item Sinh câu trả lời
	\end{enumerate}
	
	\section{Công nghệ sử dụng}
	
	\begin{itemize}
		\item Backend: FastAPI
		\item Frontend: ReactJS
		\item Vector Database: ChromaDB
		\item LLM API
	\end{itemize}
	
	% =========================
	% CHAPTER 5
	% =========================
	\chapter{Kết luận}
	
	Đề tài hướng đến xây dựng chatbot hỏi đáp tài liệu học tập bằng tiếng Việt và đánh giá hiệu quả giữa các phương pháp AI hiện đại trong bài toán hỏi đáp học liệu.
	
	% =========================
	% END DOCUMENT
	% =========================
\end{document}